\documentclass[11pt]{article}
\usepackage{ochamath}
\subjectcolor{2F5DA8}
\setsheet{線形代数・電気系院試補充}{二次形式の一般論　演習}{https://university-math-crj.pages.dev/linear-algebra/quadratic-forms-advanced/}
\namelinetrue
\begin{document}
\makesheethead{問題}
自作問題。本文の例題とは数値や設定が異なります。条件・途中式・検算を示してください。
\problem[Sylvesterの判定]
$H(t)=\begin{pmatrix}3&1&0\\1&2&1\\0&1&t\end{pmatrix}$ の正定値条件を求めよ。
\workspace{45mm}
\problem[半正定値の落とし穴]
$H=\operatorname{diag}(0,-2,1)$ は左上主小行列式が全て非負か。半正定値か。
\workspace{45mm}
\newpage
\problem[慣性]
$H=\operatorname{diag}(2,-1,0),P=\operatorname{diag}(3,2,1)$。合同変換後の固有値と慣性を求めよ。
\workspace{45mm}
\problem[制約付きRayleigh商]
$H=\operatorname{diag}(1,2,5)$、$x_1+x_2=0,\|x\|=1$ の下で $x^{\mathsf T}Hx$ の最大最小を求めよ。
\workspace{45mm}
\newpage
\problem[線形項付きの最小化]
$H=\begin{pmatrix}3&1\\1&2\end{pmatrix},b=(2,1)$ に対し $x^{\mathsf T}Hx-2b^{\mathsf T}x$ の最小点と最小値を求めよ。
\workspace{45mm}
\problem[半正定値と下に非有界]
$H=\operatorname{diag}(4,0)$ で $b=(2,0)$ と $b=(2,3)$ の最小化を比較せよ。
\workspace{45mm}
\end{document}
